\documentclass[french,11pt]{book}
\usepackage{teach}
\usepackage[french]{babel}
\usepackage{amsmath,amssymb}
\usepackage{array}
\newcommand{\R}{\mathbb{R}}
\begin{document}
\begin{center}
{\Large \textbf{Exercices -- Fonctions exponentielles}}\\[2mm]
Terminale technologique
\end{center}

\section*{Exercices}
\begin{enumerate}
  \item Une quantité augmente de 6\% par période. Écrire le coefficient multiplicateur et calculer la quantité obtenue après 5 périodes à partir de 200.
  \item Calculer $1,04^3$, $0,92^4$ et interpréter chacun des deux résultats en pourcentage d'évolution globale.
  \item Une valeur est modélisée par $f(t)=500\times 0,87^t$. Calculer $f(0)$, $f(3)$ et expliquer le sens de la base $0,87$.
  \item Déterminer à la calculatrice le plus petit entier $n$ tel que $1200\times 1,03^n>1500$.
\end{enumerate}

\end{document}
